\documentclass[conference]{IEEEtran}
\usepackage{cite}
\usepackage{amsmath,amssymb,amsfonts}
\usepackage{graphicx}
\usepackage{booktabs}
\usepackage{hyperref}
\usepackage{array}
\usepackage{xcolor}
\bibliographystyle{IEEEtran}

\begin{document}

\title{Towards Integrated AI-Assisted Legal Reasoning:\\A Survey of Indian Legal NLP Across Five Capability Axes}

\author{
\IEEEauthorblockN{Neeta Chavan, Mayur Rathod, Megh Yashwantkar, and Gajanan Barve}
\IEEEauthorblockA{
\textit{Department of Computer Science and Business Systems}\\
\textit{JSPM's Rajarshi Shahu College of Engineering, Pune, Maharashtra, India}\\
}
}

\maketitle

\begin{abstract}
India's judiciary faces a systemic crisis, with approximately fifty million cases pending across district courts, high courts, and the Supreme Court. Artificial intelligence offers a principled direction for augmenting judicial reasoning, yet existing legal AI tools remain fragmented: retrieval engines, prediction models, and explainability modules operate in isolation. This survey presents a structured review of Indian legal NLP organised along five capability axes: precedent search, case fact structuring, argument graph construction, judicial outcome simulation, and what-if analysis. For each axis, the survey catalogues representative methods, benchmark datasets, and evaluation practices. The central finding is that no existing system integrates all five capabilities into a unified pipeline. Argument graph construction and counterfactual what-if analysis remain virtually unexplored in the Indian legal context. These gaps underscore the need for future architectures that couple retrieval, structured reasoning, prediction, and interactive analysis within a single coherent framework.
\end{abstract}

%% ================================================================
%%  SECTION 1 — Introduction
%% ================================================================
\section{Introduction}

\IEEEPARstart{I}{ndia's} judicial system confronts a backlog of approximately fifty million pending cases distributed across its three-tier hierarchy of district courts, high courts, and the Supreme Court. This accumulation is not merely an administrative inconvenience; it constitutes a systemic failure of access to justice that administrative reform alone has proven insufficient to resolve. Litigants routinely wait years, sometimes decades, for resolution, eroding public confidence in the rule of law. Against this backdrop, AI-assisted legal tools represent a principled direction for augmenting, not replacing, judicial reasoning. The scale of digitised Indian legal data, including over one million Supreme Court and High Court judgments, provides a substantial empirical foundation for computational analysis~\cite{nigam-etal-2025-nyayaanumana}. These conditions make India a compelling yet underexplored testbed for legal AI research.

Despite growing research interest, current legal AI systems remain fundamentally siloed. Retrieval systems that surface relevant precedents operate independently of prediction models that forecast case outcomes, which are in turn disconnected from interactive reasoning and explainability tools. This fragmentation stands in sharp contrast to how legal reasoning actually functions: judges do not retrieve precedents in isolation, then separately apply statutes, and finally reason over facts. Rather, judicial reasoning is an integrated process wherein established facts, applicable law, and prior rulings are considered simultaneously to arrive at a disposition. Recent work has demonstrated that coupling retrieval with downstream reasoning yields substantial improvements over pipeline approaches that treat these stages independently~\cite{nigam-etal-2025-nyayarag}, and that case-based reasoning can be effectively integrated with retrieval-augmented generation to produce more legally grounded responses~\cite{wiratunga-etal-2024-cbrrag}. These findings reinforce the imperative for tighter integration across capability boundaries.

This paper presents a survey, not a system. Its contribution is twofold: first, a structured review of five capability areas in Indian legal NLP, and second, the identification of the absence of an integrated multi-capability pipeline as the field's central open problem. The five capability axes that organise this review are: (1)~precedent search, encompassing sparse and dense retrieval of relevant prior judgments; (2)~case fact structuring, covering named entity recognition, rhetorical role labelling, and statutory section tagging; (3)~argument graph construction, involving the extraction of claims, evidence, and their relational structure; (4)~judicial outcome simulation, addressing prediction of dispositions with accompanying explanations; and (5)~what-if analysis, enabling counterfactual hypothesis testing over case facts. Together, these axes span the reasoning chain from information retrieval to interactive deliberation.

The remainder of this paper is organised as follows. Section~II provides background on the Indian legal ecosystem and the specific challenges it poses for NLP. Sections~III through VII review each of the five capability axes in turn. Section~VIII presents a comparative analysis across axes, identifying integration gaps and cross-cutting themes. Section~IX discusses evaluation methodologies and metrics. Section~X outlines open problems and future research directions. Section~XI concludes the survey.

\begin{table*}[t]
\centering
\caption{Five Capability Axes in Indian Legal AI: Coverage and Identified Gaps}
\label{tab:axes}
\begin{tabular}{>{\raggedright\arraybackslash}p{3.2cm} >{\raggedright\arraybackslash}p{6.8cm} >{\raggedright\arraybackslash}p{6.8cm}}
\toprule
\textbf{Capability Axis} & \textbf{Description} & \textbf{Gap in Existing Literature} \\
\midrule
Precedent Search &
Semantic retrieval of relevant prior judgments through sparse keyword-based methods, dense vector representations, and hybrid strategies applied to Indian case law corpora. &
Retrieval systems function as standalone components; their outputs are not consumed by downstream prediction or reasoning modules, limiting their utility in end-to-end legal analysis. \\
\addlinespace
Case Fact Structuring &
Extraction and classification of structured elements from judgment text, including named entity recognition, rhetorical role labelling of sentence-level functional segments, and Indian Penal Code section tagging. &
Structured representations produced by these models are not propagated to any subsequent reasoning stage, rendering them terminal outputs rather than intermediate features in a larger pipeline. \\
\addlinespace
Argument Graph Construction &
Identification and relational mapping of argumentative components, including claims, supporting evidence, cited statutes, and invoked precedents, into directed graph representations that capture the logical structure of judicial reasoning. &
This capability is virtually absent in Indian legal AI research; existing work on legal argument mining is confined almost exclusively to Western jurisdictions such as the European Court of Human Rights. \\
\addlinespace
Judicial Outcome Simulation &
Prediction of case dispositions (acceptance or dismissal) accompanied by natural language explanations that articulate the reasoning underlying the predicted outcome. &
Prediction models operate without access to retrieval-augmented context; generated explanations remain shallow and lack grounding in specific precedents or statutory provisions. \\
\addlinespace
What-If Analysis &
Counterfactual hypothesis testing that enables users to perturb specific case facts, statutory provisions, or precedent selections and observe the resulting shifts in predicted outcomes and explanatory rationales. &
Entirely unexplored in Indian legal NLP; no existing dataset, model, or evaluation framework addresses counterfactual reasoning over Indian case law. \\
\bottomrule
\end{tabular}
\end{table*}

%% ================================================================
%%  SECTION 2 — Background: The Indian Legal Ecosystem
%% ================================================================
\section{Background: The Indian Legal Ecosystem}

India's legal system operates within the common law tradition inherited from British colonial jurisprudence, wherein judicial decisions carry precedential authority. Judicial reasoning in this framework rests on three interdependent pillars: established facts of the case, applicable statutory provisions, drawn primarily from codified instruments such as the Indian Penal Code and the Code of Civil Procedure, and binding precedents invoked through the doctrine of stare decisis. These pillars are not consulted sequentially or in isolation. A judgment emerges from their simultaneous interaction: a court assesses the facts in light of the relevant statutes, interprets those statutes in the context of authoritative prior rulings, and determines whether the factual matrix of the present case falls within the scope of established precedent. Any computational framework that aspires to model legal reasoning must therefore represent and reason over all three pillars concurrently.

Indian legal NLP presents distinct challenges that differentiate it from better-resourced Western legal domains. First, Indian court judgments are frequently multilingual: a single document may interleave English with Hindi, Urdu, or regional languages without explicit demarcation or translation, complicating tokenisation and representation learning. Second, there is no standardised formatting convention across Indian courts or time periods; judgments from different benches, courts, and decades exhibit wide variation in structure, citation style, and rhetorical organisation, impeding rule-based and template-driven extraction methods. Third, the availability of high-quality annotated data remains limited relative to established benchmarks in Western jurisdictions such as the European Court of Human Rights or United States federal case law. The LegalEval shared task highlighted the difficulty of constructing consistent annotations for Indian legal text, revealing substantial inter-annotator disagreement on tasks such as rhetorical role labelling~\cite{modi-etal-2023-semeval}. Domain-adapted language models such as InLegalBERT have demonstrated that general-purpose pre-trained representations transfer poorly to this domain, necessitating targeted pre-training on Indian legal corpora~\cite{paul-etal-2023-inlegalbert}.

Despite these challenges, Indian legal NLP has undergone rapid growth in recent years, catalysed by the creation of several large-scale datasets. The Indian Legal Documents Corpus (ILDC) introduced a benchmark of over thirty-five thousand Supreme Court cases annotated for judgment prediction and explanation~\cite{malik-etal-2021-ildc}. PredEx extended this line of work by providing expert-annotated explanations grounded in specific textual spans of the judgment~\cite{nigam-etal-2024-predex}. Most recently, NyayaAnumana established the largest Indian legal judgment prediction corpus to date, spanning multiple court levels and accompanied by a domain-specialised language model~\cite{nigam-etal-2025-nyayaanumana}. These resources have enabled measurable progress on individual tasks such as binary judgment prediction and rhetorical role classification. However, no existing work has unified these capabilities under a coherent multi-capability architecture. Each dataset and model addresses a single axis in isolation, leaving the integration of retrieval, structuring, reasoning, prediction, and interactive analysis as the field's principal unresolved challenge. This gap motivates the five-axis taxonomy that structures the remainder of this survey.


%% ================================================================
%%  SECTION 3 — Precedent Search
%% ================================================================
\section{Precedent Search}

Precedent retrieval constitutes the foundational layer of any legal AI pipeline. Within the common law tradition, prior judicial rulings establish binding interpretations of statutory provisions and factual patterns, and identifying the most pertinent precedents is the first cognitive step in constructing a legal argument. Recent work has demonstrated that grounding downstream generation in retrieved case law yields substantial improvements over approaches that rely solely on parametric knowledge~\cite{nigam-etal-2025-nyayarag}, and that coupling retrieval with structured case-based reasoning further enhances the legal fidelity of generated outputs~\cite{wiratunga-etal-2024-cbrrag}. Retrieval quality therefore directly constrains the quality of every subsequent task in the reasoning chain, from fact analysis to outcome prediction.

\subsection{Keyword-Based and Sparse Retrieval}

Keyword-based search has historically served as the primary mechanism for precedent discovery in Indian legal practice. Indian Kanoon, the most widely used public legal search platform, operates fundamentally on term-matching principles. The standard formalisation of sparse retrieval is the BM25 scoring function, which ranks documents by term frequency modulated by inverse document frequency across the corpus. Despite the proliferation of neural retrieval methods, BM25 remains a remarkably strong baseline; empirical evaluations on legal case retrieval benchmarks have shown that it matches or exceeds several neural approaches, particularly when queries contain distinctive legal terminology~\cite{rosa-etal-2021-bm25}. The principal limitation of sparse retrieval is vocabulary mismatch: a query and a semantically relevant document may share no surface-level terms. For example, a query referencing ``breach of fiduciary duty'' may fail to retrieve a judgment that discusses the same legal concept using the phrase ``violation of trust obligations.''

\subsection{Dense Retrieval}

Dense retrieval addresses the vocabulary mismatch problem by encoding queries and documents into continuous vector representations, where retrieval is performed via cosine or dot-product similarity in the embedding space. Transformer-based encoders learn to map semantically related texts to proximal regions of the vector space, capturing conceptual overlap that keyword methods cannot detect. In the legal domain, dense retrieval has been applied to the COLIEE competition's case law retrieval task, where dual-encoder architectures combined with summarisation-based re-ranking have demonstrated competitive performance against sparse baselines~\cite{althammer-etal-2021-dossier}. Scaling dense retrieval to large corpora requires approximate nearest-neighbour indexing; FAISS provides the standard infrastructure for this purpose, enabling sub-linear search over billion-scale vector collections on commodity GPU hardware~\cite{johnson-etal-2021-faiss}. The trade-off is computational: dense retrieval demands substantial resources for both index construction and query-time inference relative to the lightweight inverted indices used by BM25.

\subsection{Domain-Adapted Models for Indian Legal Retrieval}

General-purpose language models pre-trained on news articles or encyclopaedic text transfer poorly to legal corpora, where specialised vocabulary, complex citation conventions, and unusually long sentences diverge markedly from the training distribution. This motivates domain adaptation, that is, continued pre-training of transformer encoders on in-domain legal text. InLegalBERT addresses this gap for the Indian context by pre-training a BERT-based model on a large corpus of Supreme Court and High Court judgments~\cite{paul-etal-2023-inlegalbert}. Evaluations demonstrate that InLegalBERT consistently outperforms both generic BERT and multilingual variants on Indian legal retrieval and classification benchmarks, confirming that domain-specific distributional patterns are essential for effective legal text representation.

\subsection{Retrieval-Augmented Generation for Legal Reasoning}

Retrieval-augmented generation (RAG) represents the current frontier in legal information retrieval by coupling a retrieval module with a generative large language model. Rather than relying exclusively on parametric knowledge, the generative component receives retrieved documents as context, grounding its output in actual precedents. NyayaRAG applies this paradigm specifically to Indian legal judgment prediction, demonstrating that retrieval-grounded generation substantially outperforms standalone generation across multiple court levels~\cite{nigam-etal-2025-nyayarag}. CBR-RAG extends the paradigm by incorporating case-based reasoning, retrieving cases that are structurally analogous rather than merely semantically proximate, thereby improving the legal relevance of the retrieved context~\cite{wiratunga-etal-2024-cbrrag}. Despite this progress, a critical gap persists: retrieval outputs in existing Indian legal AI systems are consumed by generation modules alone and are not connected to structured fact analysis, argument construction, or outcome prediction in any unified pipeline.

%% ================================================================
%%  SECTION 4 — Case Fact Structuring
%% ================================================================
\section{Case Fact Structuring}

Indian court judgments are lengthy, unstructured documents, frequently spanning tens of thousands of words, that interleave procedural history, factual narratives, statutory interpretation, counsels' arguments, and the operative ruling without standardised demarcation. Before any downstream reasoning system can operate effectively, this raw text must be decomposed into its functional constituents. This section surveys the three primary structuring tasks that have been addressed in Indian legal NLP: named entity recognition, rhetorical role classification, and statutory section tagging.

\subsection{Legal Named Entity Recognition}

Named entity recognition in the Indian legal domain involves the identification of entity types that are specific to judicial text and substantially different from the person--location--organisation categories of general-purpose NER. Kalamkar et al.\ define a fine-grained schema of fourteen entity types tailored to Indian court judgments, including court, judge, petitioner, respondent, lawyer, date, case number, statute, provision, precedent, witness, other person, organisation, and geographical entity~\cite{kalamkar-etal-2022-ner}. Annotations were constructed over Supreme Court and High Court judgments, and transformer-based sequence labelling models were evaluated against this benchmark. Fine-grained legal NER enables critical downstream capabilities: precedent entities can be linked to their source judgments for retrieval, party entities support coreference resolution across the document, and statute entities provide a bridge to the statutory knowledge base. Without this structuring step, downstream modules must operate over raw text with no access to the relational scaffolding that entities provide.

\subsection{Rhetorical Role Classification}

Rhetorical role labelling assigns each sentence in a judgment a functional category indicating its role in the document's argumentative structure, such as whether it states a fact found by the court, recites a ruling from a prior case, articulates a legal issue under consideration, presents the court's reasoning, or delivers the final decision. The SemEval-2023 Task~6 (LegalEval) established a multi-court benchmark for this task across Indian jurisdictions, revealing both the feasibility and the difficulty of automated rhetorical segmentation~\cite{modi-etal-2023-semeval}. A central challenge is that the rhetorical role of a sentence frequently depends on long-range context: a sentence stating ``the appeal is dismissed'' carries different rhetorical significance depending on whether it appears within a recitation of a prior ruling or as the operative order of the present case. Hierarchical and sequential models that propagate contextual information across sentence boundaries have been employed to address this dependency. Rhetorical segmentation is what permits an AI system to isolate, for example, the ratio decidendi of a judgment rather than treating the document as an undifferentiated sequence.

\subsection{IPC Section Tagging and Statute Linking}

The task of statute identification involves automatically determining which sections of the Indian Penal Code or other legislative instruments are invoked, discussed, or applied within a given judgment. LegalEval includes statute identification as one of its three constituent subtasks, providing annotated data and evaluation protocols for this problem~\cite{modi-etal-2023-semeval}. Statute tagging is particularly consequential because it connects the factual matrix of a case to the normative legal framework, a prerequisite for any downstream argument construction or outcome prediction module. A persistent challenge is implicit reference: judgments frequently discuss the substance of a statutory provision, including its elements, its judicial interpretation, and its applicability to the facts, without explicitly naming the section number, requiring models to perform semantic matching between case text and statutory language. Taken together, NER, rhetorical role classification, and statute tagging address complementary dimensions of document structure. However, these three tasks have been studied largely in isolation within the Indian legal NLP literature; no existing system propagates all three structured outputs into a unified downstream reasoning pipeline, leaving the structuring layer as a set of terminal outputs rather than intermediate representations in a larger architecture.

%% ================================================================
%%  SECTION 5 — Argument Graph Construction
%% ================================================================
\section{Argument Graph Construction}

Argument graph construction is the most structurally novel of the five capability axes examined in this survey. Legal reasoning is fundamentally argumentative: a judicial decision emerges from the interplay of competing claims advanced by opposing parties, supporting evidence drawn from the factual record, statutory provisions that supply normative authority, and precedential rulings invoked by analogy or distinction. Representing these components and their logical relations as a directed graph renders the inferential skeleton of a case both explicit and machine-readable. This capability area has the least direct prior work in the Indian legal AI context; the theoretical foundations and empirical benchmarks must therefore be drawn from the broader argument mining literature and from research conducted on Western legal corpora.

\subsection{Argument Mining: Foundations}

Argument mining encompasses the automatic identification of argumentative components, including claims, premises, and evidence, together with the relations that hold between them, such as support, attack, and rebuttal. In the legal domain, Habernal et al.\ conduct the most comprehensive study to date, mining arguments from decisions of the European Court of Human Rights using a taxonomy of seventeen argumentative component types~\cite{habernal-etal-2023-mining}. Their work demonstrates that transformer-based models can recover argumentative structure from long legal documents with reasonable accuracy, though performance degrades as document length and argumentative complexity increase. The field has progressively shifted from sentence-level classification, in which individual sentences are labelled as claims or premises, to graph-level modelling, where nodes represent argumentative units and directed edges encode support, contradiction, or evidential relations between them. This graph-level perspective is essential for legal applications, where a single precedent may simultaneously support one claim and undermine another.

\subsection{Argument Structure in Legal Corpora}

The application of argument mining to legal text has concentrated predominantly on European and American jurisdictions. LexGLUE provides a benchmark suite for legal language understanding that encompasses multiple tasks requiring comprehension of argumentative structure, evaluated primarily on ECHR and other Western legal corpora~\cite{chalkidis-etal-2022-lexglue}. Indian courts are absent from this benchmark, establishing a direct empirical gap. The canonical structure of a legal argument graph maps naturally onto the components of a judgment: claims correspond to the legal issues in dispute, evidence nodes capture the facts established at trial, statute nodes supply the normative warrant grounding the legal analysis, and precedent nodes furnish analogical support or distinguishing authority. Once such a graph is constructed, weak-point detection, defined as the identification of claims insufficiently supported by evidence or statutory authority, follows directly from graph completeness analysis, offering a principled basis for evaluating the strength of competing legal positions.

\subsection{The Indian Legal Gap and Graph Database Representations}

Argument graph construction for Indian legal judgments remains an open problem: as of the time of writing, no published dataset, annotation schema, or trained model addresses this task in the Indian context. Several factors account for this gap. Indian judgments are substantially longer and more discursive than the ECHR decisions on which existing argument mining models have been developed. The argumentative structure of Indian judgments is less formulaic, with reasoning frequently distributed across non-contiguous sections of the document. Most critically, the annotated data necessary to train and evaluate argument mining models for Indian courts does not exist. From a representational standpoint, graph databases such as Neo4j provide a natural storage and query substrate for legal argument structures, with nodes instantiated for claims, evidence, statutes, and precedents, and typed edges encoding support, citation, contradiction, and analogy relations. The outputs of the preceding sections, namely retrieved precedents from Section~III, and NER entities, rhetorical segments, and statute tags from Section~IV, constitute the natural inputs to argument graph construction, positioning this capability as the connective layer within an integrated pipeline architecture. Indian legal argument mining therefore stands as a high-priority and largely uncharted research direction.


%% ================================================================
%%  SECTION 6 — Judicial Outcome Simulation
%% ================================================================
\section{Judicial Outcome Simulation}

Judicial outcome simulation is the most extensively studied of the five capability axes surveyed in this paper, with a research lineage that extends from classical machine learning classifiers through transformer-based models to contemporary large language models. The task can be defined precisely: given the text of a pending or historical case, predict the binary or multi-class outcome (accepted, rejected, or partially accepted) and, where applicable, generate a natural language explanation articulating the reasoning behind that prediction. The field is presently at a transition point in how this task is conceptualised and evaluated~\cite{nigam-etal-2024-rethinking}. Prediction accuracy alone is insufficient for legal contexts, where the reasoning process underlying a decision carries at least as much weight as the dispositional outcome itself.

\subsection{Classical and Early Neural Approaches}

The earliest approaches to legal judgment prediction relied on hand-crafted features, including bag-of-words representations, TF-IDF vectors, and manually defined statutory indicators, fed into classical classifiers such as support vector machines and logistic regression. These models achieved modest accuracy on small, curated datasets but exhibited poor generalisation across court levels, case types, and temporal periods. The subsequent transition to convolutional and recurrent neural networks enabled models to operate on raw text rather than engineered feature sets, eliminating the manual feature design bottleneck. However, these architectures introduced new challenges around long document handling, as Indian court judgments routinely exceed the input capacity of standard recurrent models. The introduction of LegalBERT represented an explicit response to the inadequacy of general-purpose neural models for legal text, demonstrating that domain-specific pre-training on legal corpora yields representations better suited to the distributional characteristics of judicial language~\cite{chalkidis-etal-2020-legalbert}.

\subsection{Transformer-Based Prediction}

The adoption of transformer architectures marked a turning point for legal judgment prediction. Pre-trained models such as LegalBERT~\cite{chalkidis-etal-2020-legalbert} and InLegalBERT~\cite{paul-etal-2023-inlegalbert} delivered substantial accuracy gains by leveraging domain-specific pre-training on large legal corpora. In the Indian context, the Indian Legal Documents Corpus (ILDC) made systematic transformer evaluation possible by providing over 35,000 Supreme Court cases annotated with binary outcome labels and a multi-sentence explanation layer~\cite{malik-etal-2021-ildc}. A persistent architectural challenge is document length: Indian judgments routinely exceed the 512-token context window of standard BERT models, necessitating hierarchical encoders, overlapping chunk strategies, or Longformer-style architectures that extend attention to thousands of tokens. These adaptations introduce additional complexity in training and inference while still imposing an upper bound on the amount of judgment text that can inform a single prediction.

\subsection{Explainability as a First-Class Requirement}

Explainability in legal outcome prediction is not an optional post-hoc addition but a functional requirement. A prediction system that outputs only a label provides no actionable guidance to lawyers or judges and cannot be meaningfully audited for bias, error, or unfairness. PredEx addresses this gap by introducing a corpus of over 15,000 Indian court cases annotated with natural language explanations aligned to prediction outcomes~\cite{nigam-etal-2024-predex}. These explanations are not post-hoc rationalisations generated by the model but are composed by legal professionals to reflect the actual reasoning underlying each judgment. This methodology represents a significant advance over attention-based explanation proxies, which prior research has shown to be unreliable indicators of model reasoning. By grounding explainability in expert annotation rather than model introspection, PredEx establishes a standard against which future explanation systems can be measured.

\subsection{Large Language Models and the Scale Frontier}

The current state of the art in judicial outcome simulation is defined by large language models trained or fine-tuned on legal corpora at unprecedented scale. NyayaAnumana provides over 700,000 Indian court case proceedings spanning district courts, high courts, and the Supreme Court, accompanied by INLegalLlama, a domain-specialised variant of LLaMA evaluated on this corpus~\cite{nigam-etal-2025-nyayaanumana}. Models trained on corpora of this scale begin to internalise the distributional patterns of judicial reasoning across different court levels, case categories, and statutory domains. However, a critical evaluation challenge has emerged: standard legal judgment prediction benchmarks overestimate real-world performance when cases are selected retrospectively, as benchmarks disproportionately include cases with clear outcomes rather than genuinely contested disputes~\cite{nigam-etal-2024-rethinking}. Retrieval-augmented generation represents the current best-performing configuration, as retrieved precedents provide contextual grounding that parametric knowledge alone cannot supply~\cite{nigam-etal-2025-nyayarag}. This finding reinforces the dependency of prediction quality on retrieval quality, further motivating the integration of these capabilities.

\subsection{Remaining Limitations}

Three concrete limitations characterise the current state of judicial outcome simulation. First, prediction models are trained and evaluated in isolation from the retrieval and structuring systems surveyed in Sections~III and~IV, meaning that the quality of retrieved precedents and structured case facts has no pathway to influence prediction accuracy. Second, evaluation is overwhelmingly concentrated on Supreme Court data. District courts and high courts, which collectively handle the vast majority of India's pending caseload, remain understudied despite exhibiting different linguistic patterns, case complexities, and outcome distributions. Third, long-document handling remains unsolved at scale: even modern LLMs with extended context windows have not been rigorously evaluated on Indian judgments that exceed 50,000 words. More broadly, the legal judgment prediction community has been subject to the critique that evaluation metrics are frequently selected without rigorous justification, with accuracy dominating at the expense of more informative measures such as macro-averaged F1, calibration error, and per-class recall~\cite{medvedeva-etal-2020-echr}.

%% ================================================================
%%  SECTION 7 — What-If Analysis
%% ================================================================
\section{What-If Analysis}

What-if analysis is the least explored of the five capability axes surveyed in this paper and the one most directly oriented toward the practising lawyer rather than the academic researcher. The core task is counterfactual reasoning over case facts: given a case and a predicted outcome, identify what change to the facts, statutory framing, or precedent basis would alter that outcome. This is distinct from prediction, which asks what the outcome is, and from explanation, which asks why a particular prediction was made. What-if analysis asks what would need to be different. Within the broader literature on explainable artificial intelligence, counterfactual explanation is recognised as one of the most practically actionable forms of model interpretability, valued precisely because it translates model behaviour into concrete, decision-relevant guidance~\cite{mersha-etal-2024-xai}.

\subsection{Counterfactual Explanation in NLP}

Counterfactual explanation in natural language processing involves generating minimally modified inputs that produce a different model output. In a text classification setting, a counterfactual explanation identifies which words, phrases, or factual assertions, if altered or removed, would flip the predicted label. The general XAI literature provides a taxonomy of methods that includes gradient-based saliency maps, attention weight analysis, and perturbation-based approaches, each offering a different lens on the relationship between input features and model decisions~\cite{mersha-etal-2024-xai}. A key distinction separates counterfactual explanation from feature importance. Feature importance identifies which input elements contributed most to the current prediction, answering the question of what mattered. Counterfactual explanation identifies what would need to change to produce a different prediction, answering the question of what could be different. For a lawyer constructing or challenging a legal argument, the latter formulation is considerably more useful, as it directly maps onto the strategic decisions involved in case preparation: which facts to emphasise, which statutory provisions to invoke, and which precedents to distinguish.

\subsection{XAI in Legal Decision-Making}

The application of explainability methods to legal AI systems raises both technical and normative considerations. From a normative perspective, explainability in automated legal decision-making is not merely a technical desideratum but a procedural requirement in jurisdictions with transparency obligations, where affected parties have a right to understand the basis of decisions that affect their legal interests~\cite{deakin-etal-2023-xai-law}. This creates a fundamental tension: the highest-performing models in legal prediction are typically the most opaque, while the models most amenable to explanation tend to sacrifice predictive accuracy. Local explanation methods such as LIME and additive feature attribution approaches such as SHAP have been applied to legal classifiers, but these methods are designed for feature importance rather than the counterfactual use case that legal practitioners require. PredEx's expert-annotated explanations represent progress toward legally meaningful interpretability, though they address the explanation of observed predictions rather than the exploration of hypothetical alternative scenarios~\cite{nigam-etal-2024-predex}. Bridging this gap, from explaining what happened to exploring what could happen, remains an open challenge.

\subsection{The Indian Legal Gap and Practical Motivation}

Counterfactual what-if analysis has not been applied to Indian legal AI as of the time of writing. The closest analogues are sensitivity analyses conducted on European legal prediction models, and these have not been adapted to the institutional, linguistic, or statutory context of the Indian judiciary. The practical value of such a capability for the Indian legal system is substantial. A lawyer could test whether changing the section of the Indian Penal Code under which a charge is framed would alter the predicted outcome. A judge could explore whether the addition of a particular precedent to the case record would shift the balance of the analysis. Such hypothesis testing is particularly valuable given the resource constraints and data scarcity challenges that characterise legal AI development in jurisdictions outside of Western Europe and North America~\cite{medvedeva-etal-2020-echr}. What-if analysis permits structured exploration of the decision space without requiring model retraining or additional annotated data, making it a low-cost complement to the prediction and argument graph capabilities surveyed in Sections~V and~VI. This positions counterfactual reasoning as a high-priority direction for future Indian legal NLP research, one that would close the loop between retrieval, structured analysis, prediction, and interactive deliberation~\cite{mersha-etal-2024-xai, deakin-etal-2023-xai-law}.


%% ================================================================
%%  SECTION 8 — Datasets
%% ================================================================
\section{Datasets}

The availability and quality of annotated datasets has been the primary bottleneck in Indian legal NLP. Unlike Western legal AI, which benefits from decades of structured case law repositories such as LexisNexis and PACER, Indian legal NLP has had to construct its own benchmarks largely from scratch, frequently relying on manual annotation by legal professionals. The seven datasets surveyed in this section represent the field's principal resources as of 2025, and they vary substantially in size, court coverage, annotation type, and public availability.

\begin{table*}[t]
\centering
\caption{Primary Datasets for Indian Legal NLP Research (2021--2025)}
\label{tab:datasets}
\begin{tabular}{>{\raggedright\arraybackslash}p{2.4cm} >{\centering\arraybackslash}p{1.8cm} >{\raggedright\arraybackslash}p{1.8cm} >{\raggedright\arraybackslash}p{2.4cm} >{\raggedright\arraybackslash}p{4.2cm} >{\centering\arraybackslash}p{1.2cm}}
\toprule
\textbf{Dataset} & \textbf{Reference} & \textbf{Size} & \textbf{Court Coverage} & \textbf{Annotation Type} & \textbf{Avail.} \\
\midrule
ILDC & \cite{malik-etal-2021-ildc} & 35,000 cases & Supreme Court of India & Binary outcome labels, multi-sentence explanations & Public \\
\addlinespace
PredEx & \cite{nigam-etal-2024-predex} & 15,000+ cases & Supreme Court of India & Expert-annotated natural language explanations aligned to outcomes & Public \\
\addlinespace
NyayaAnumana & \cite{nigam-etal-2025-nyayaanumana} & 700,000+ cases & Supreme Court, High Courts, District Courts & Outcome labels, court-level metadata & Public \\
\addlinespace
NyayaRAG Corpus & \cite{nigam-etal-2025-nyayarag} & Subset of Indian judgments & Supreme Court and High Courts & Retrieval-annotated with relevant precedent links & Partial \\
\addlinespace
LegalEval (SemEval-2023 Task 6) & \cite{modi-etal-2023-semeval} & 9,380 sent. (NER); 14,864 sent. (roles) & Supreme Court and High Courts & NER labels (14 types), rhetorical role labels, statute tags & Public \\
\addlinespace
InLegalBERT Corpus & \cite{paul-etal-2023-inlegalbert} & 5.4M sentences & Supreme Court and High Courts & Pre-training corpus, no task-specific labels & Partial \\
\addlinespace
Indian Kanoon & -- & Millions of documents & All court levels & No annotation, raw text only & Public \\
\bottomrule
\end{tabular}
\end{table*}

The progression in dataset scale from ILDC's 35,000 cases to NyayaAnumana's 700,000 cases represents a step change that enables LLM-scale fine-tuning for the first time in the Indian legal context~\cite{nigam-etal-2025-nyayaanumana}. Scale alone, however, does not compensate for annotation depth. NyayaAnumana provides outcome labels and court-level metadata but lacks the expert-annotated natural language explanations that PredEx supplies at smaller scale~\cite{nigam-etal-2024-predex}. District court data remains severely underrepresented across all annotated datasets relative to the proportion of India's pending caseload that these courts handle in practice. The InLegalBERT corpus and Indian Kanoon provide raw text at large scale but contain no task-specific annotations, limiting their utility to pre-training and unsupervised analysis.

Three annotation gaps persist that no existing dataset addresses. First, no dataset provides argument-level annotations that link claims, evidence, statutes, and precedents within a single judgment, which renders argument graph construction an unsupervised or weakly supervised problem by necessity. Second, no dataset covers multilingual judgments in which Hindi or regional language passages are embedded within English text, despite the prevalence of such code-mixing in Indian judicial documents. Third, no dataset provides counterfactual annotations of the kind required to train or evaluate a what-if analysis system. These three gaps correspond directly to the open research directions identified in Sections~V and~VII, and their resolution will require coordinated annotation efforts involving both NLP researchers and legal domain experts~\cite{modi-etal-2023-semeval}.

%% ================================================================
%%  SECTION 9 — Evaluation Metrics
%% ================================================================
\section{Evaluation Metrics}

Evaluation methodology is as consequential as model architecture in legal AI research, yet it has received comparatively little critical attention. The choice of metric determines what a system is optimised for, and in a legal context, optimising for the wrong objective can produce systems that appear performant on benchmarks but fail systematically in deployment. The majority of legal AI studies rely on accuracy as their sole evaluation measure, a choice that is rarely justified and that masks systematic failure modes on minority outcome classes~\cite{medvedeva-etal-2020-echr}. This section surveys the metrics employed across the retrieval, classification, and generation tasks covered by this survey.

\subsection{Retrieval Metrics}

Mean Average Precision (MAP) is the standard metric for evaluating ranked retrieval, capturing both the precision of retrieved documents and the quality of their ranking order. MAP computes the average precision at each relevant document in the ranked list and then averages across queries, yielding a single score that reflects both recall and rank position. Normalised Discounted Cumulative Gain (nDCG) provides a complementary perspective by applying a logarithmic discount to documents at lower ranks, placing greater weight on highly ranked relevant results. This weighting reflects the behaviour of legal practitioners, who typically examine only the top results of a retrieval query. Both metrics require relevance judgments that are expensive to collect in legal domains, as the relevance of a prior case to a present dispute is a matter of legal interpretation rather than surface similarity. BM25 evaluation on legal case retrieval benchmarks~\cite{rosa-etal-2021-bm25} and dense retrieval experiments in the COLIEE competition~\cite{althammer-etal-2021-dossier} employ both MAP and nDCG as their primary evaluation measures.

\subsection{Classification, Generation, and Explainability Metrics}

For binary and multi-class outcome prediction, accuracy and macro-averaged F1 are the standard measures. Macro-F1 is preferable to accuracy when class distributions are imbalanced, a condition that is common in Indian legal datasets where one outcome class (typically dismissal) dominates. Relying solely on accuracy in such settings can produce misleadingly high scores, as a trivial majority-class classifier achieves high accuracy without learning any discriminative signal~\cite{medvedeva-etal-2020-echr}. For named entity recognition, per-entity-type F1 is the appropriate measure, since aggregate F1 can mask poor performance on rare but legally significant entity types such as witness names and provision numbers~\cite{kalamkar-etal-2022-ner}. For the evaluation of generated explanations, automated metrics such as ROUGE and BERTScore measure surface overlap with reference texts but cannot assess legal soundness. Expert Likert scale assessments, in which legal professionals rate the adequacy, accuracy, and completeness of generated explanations on a structured scale, currently represent the most reliable evaluation methodology. This approach has been employed in the evaluation of both INLegalLlama~\cite{nigam-etal-2025-nyayaanumana} and PredEx~\cite{nigam-etal-2024-predex}. Human evaluation of this kind is resource-intensive but remains the only dependable measure of whether a generated explanation is legally sound rather than merely fluent.


%% ================================================================
%%  SECTION 10 — Gaps and Contribution
%% ================================================================
\section{Gaps and Contribution}

The central finding of this survey is that no existing Indian legal AI system integrates precedent retrieval, case fact structuring, argument graph construction, outcome prediction, and what-if analysis into a unified pipeline. Each of these five capabilities has been studied in isolation, and the connections that would link them into a coherent reasoning chain remain unimplemented in any published system. Retrieved precedents are not fed into structured fact analysis modules. Structured facts and entity annotations are not consumed by argument graph constructors. Argument graphs do not inform prediction models. Prediction outputs do not serve as the basis for counterfactual exploration. The closest existing work to pipeline integration couples retrieval with generation~\cite{nigam-etal-2025-nyayarag} or retrieval with case-based reasoning~\cite{wiratunga-etal-2024-cbrrag}, but even these systems leave structuring, argument modelling, and what-if analysis entirely outside their scope.

Three task-level gaps stand out as particularly consequential. First, argument graph construction for Indian legal judgments has no published dataset, trained model, or evaluation benchmark, making it the field's most significant structural blind spot. Second, counterfactual what-if analysis has not been attempted in the Indian legal context, despite being the capability most directly useful to practising lawyers preparing or challenging legal arguments. Third, evaluation across the surveyed tasks is systematically skewed toward Supreme Court data~\cite{malik-etal-2021-ildc, nigam-etal-2024-predex}, leaving district and high court levels without reliable benchmarks despite their collective responsibility for the overwhelming majority of India's pending cases. The annotation scope of existing shared tasks such as LegalEval further reflects this imbalance~\cite{modi-etal-2023-semeval}.

Three data-level gaps compound these task-level deficiencies. First, no annotated dataset captures argument-level structure within Indian judgments, linking claims to their supporting evidence, statutory warrants, and cited precedents. Second, no dataset addresses the multilingual character of Indian court proceedings, where Hindi and regional language passages embedded within English judgments remain unprocessed by all existing models. Third, the large-scale datasets that enable LLM fine-tuning, such as NyayaAnumana~\cite{nigam-etal-2025-nyayaanumana}, do not provide the expert-level explanation annotations available in smaller datasets such as PredEx~\cite{nigam-etal-2024-predex}. No existing resource combines both scale and annotation depth simultaneously.

The contribution of this survey is precisely circumscribed. It provides the first structured review of Indian legal NLP organised around five functional capability axes, maps the existing literature to each axis, identifies the datasets and evaluation methods associated with each, and produces a consolidated synthesis of what remains unsolved. The primary finding is that the field has matured sufficiently in individual capabilities to make integration a realistic next step, and that the absence of such integration is now the binding constraint on further progress. This survey serves as a reference point for researchers and practitioners seeking to build toward a unified Indian legal AI architecture.

%% ================================================================
%%  SECTION 11 — Ethical Considerations
%% ================================================================
\section{Ethical Considerations}

Any AI system trained on historical Indian court judgments inherits the biases embedded in those judgments. Indian legal history contains well-documented disparities in outcomes correlated with gender, caste, economic status, and geographic region. A prediction model trained to replicate historical outcome distributions will replicate historical disparities unless active debiasing measures are applied during training and evaluation. This risk is not unique to India; bias amplification has been identified as a general problem in legal prediction models across jurisdictions~\cite{medvedeva-etal-2020-echr}. The risk is, however, compounded in any system that influences real legal decisions, because model outputs may be treated as authoritative even when they reflect statistical artefacts rather than legal reasoning. Bias auditing across demographic and geographic subgroups must therefore be a mandatory component of any deployed Indian legal AI system, not an optional post-deployment analysis.

In legal contexts, explainability is not a technical convenience but a procedural requirement. Parties to litigation have a right to understand the basis on which decisions affecting their legal interests are made, and this right extends to any AI-assisted component of the decision-making process~\cite{deakin-etal-2023-xai-law}. A specific tension arises between the opacity of high-performing large language models and the transparency obligations that govern judicial proceedings. Generated text that is fluent and well-structured may nonetheless be legally unsound, and the distinction between fluency and soundness is not detectable by automated metrics alone. The move toward expert-annotated explanations represented by PredEx~\cite{nigam-etal-2024-predex} is the correct methodological direction, as it grounds the evaluation of explanations in domain expertise rather than surface-level text similarity. Fluency of generated text must not be mistaken for legal soundness.

A clear and unambiguous distinction must be drawn between a legal AI system as a decision-support tool and one as an autonomous adjudicator. The systems surveyed in this paper are appropriately characterised as the former: tools that organise information, surface relevant precedents, structure arguments, and estimate likely outcomes to assist human legal professionals in their reasoning. They are not, and should not be positioned as, replacements for judicial reasoning. Even the most capable current systems operate under significant limitations in coverage, robustness, and interpretability that preclude autonomous deployment~\cite{nigam-etal-2025-nyayaanumana, mersha-etal-2024-xai}. The responsible development of Indian legal AI requires sustained collaboration among NLP researchers, legal professionals, and policymakers, ensuring that technical capability is matched by institutional accountability and that the augmentation of legal reasoning does not erode the procedural safeguards on which the rule of law depends.

\bibliography{survey_references}
\end{document}

